\documentclass[11pt,a4paper]{article}

\usepackage[utf8]{inputenc}
\usepackage[margin=2.5cm]{geometry}
\usepackage{amsmath,amssymb,amsthm}
\usepackage{mathtools}
\usepackage{booktabs}
\usepackage{hyperref}
\usepackage{microtype}

\hypersetup{
    colorlinks=true,
    linkcolor=blue,
    citecolor=red,
    urlcolor=blue
}

% Theorem Environments
\newtheorem{theorem}{Theorem}[section]
\newtheorem{lemma}[theorem]{Lemma}
\newtheorem{proposition}[theorem]{Proposition}
\newtheorem{corollary}[theorem]{Corollary}
\theoremstyle{definition}
\newtheorem{definition}[theorem]{Definition}
\newtheorem{remark}[theorem]{Remark}
\newtheorem{example}[theorem]{Example}

\newcommand{\F}{\mathbb{F}_2}
\newcommand{\Z}{\mathbb{Z}}
\newcommand{\V}{\mathbb{V}}
\newcommand{\wt}{\mathrm{wt}}
\newcommand{\rad}{\mathrm{rad}}
\newcommand{\supp}{\mathrm{supp}}
\newcommand{\rank}{\mathrm{rank}}

\title{\textbf{Odd-Order Reduction and a Divisibility Obstruction for Cubic Homogeneous Rotation-Symmetric Bent Functions}}

\author{
  \textbf{Valer et al.}\thanks{Working paper and independent research dossier on cryptographic Boolean functions.}
}

\date{October 2026}

\begin{document}

\maketitle

\begin{abstract}
The St\u{a}nic\u{a}--Maitra conjecture (2008) asserts that there exist no homogeneous rotation-symmetric bent Boolean functions (HRSBF) of algebraic degree $d > 2$. While Cusick and Sanger (2017) proved the non-existence of cubic HRSBFs for dimensions $n = 2p$ ($p$ an odd prime), the conjecture has remained open for general composite dimensions and powers of two. In this paper, we establish an analytical odd-order reduction framework: for any linear automorphism $T$ of odd order $m$ acting on $\F^n$, a quadratic Boolean function $q$ is balanced on $\F^n$ if and only if its restriction to the fixed subspace $U = \mathrm{Fix}(T)$ is balanced on $U$. Consequently, any cubic bent function on $\F^n$ invariant under $T$ restricts to a bent function on $U$. Applying this to the cyclic shift $T = \sigma^t$ for $n = t \cdot m$ ($t = 2^s$, $m$ odd), we prove that all short orbits of length $l < t$ cancel modulo 2, so that the restriction is spanned strictly by complete orbits of full length $t$. By analyzing the resulting reduced spaces, we prove the non-existence of cubic HRSBFs for $t \in \{2, 4, 8\}$. Furthermore, for $t = 16$, we prove that the evaluation matrix of the 42 complete orbits generates a 16-divisible binary linear code; via an exact weight polarization expansion within the divisible codes framework of Ward~\cite{ward1981,ward2001}, every codeword satisfies $\wt(g) \equiv 0 \pmod{256}$, forcing the Walsh transform at zero to satisfy $W_g(0) \equiv 0 \pmod{512}$, which precludes bentness since $|W_g(0)| = 256$. Combining these reductions, we prove that no cubic HRSBF exists in $n$ variables for any even $n$ not divisible by 32, resolving the cubic case for the infinite families $n \in \{2m, 4m, 8m, 16m\}$ for all odd $m$. Finally, we report the complete computational exclusion of quartic HRSBFs for $n = 8$ and $n = 10$.
\end{abstract}

\noindent\textbf{Keywords:} Cryptographic Boolean functions, Bent functions, Rotation-symmetric functions, St\u{a}nic\u{a}--Maitra conjecture, Odd-order reduction, Divisible codes.

%======================================================================
\section{Introduction}
\label{sec:intro}
%======================================================================

Bent Boolean functions, introduced by Rothaus~\cite{rothaus1976}, achieve the maximum possible nonlinearity in even dimensions and play a central role in stream ciphers, S-box design, and coding theory. Rotation-symmetric Boolean functions (RSBFs), whose algebraic normal forms (ANF) are invariant under cyclic shifts of the input coordinates, were introduced by Pieprzyk and Qu~\cite{pieprzyk1999} and further investigated by St\u{a}nic\u{a} and Maitra~\cite{stanica2008} due to their hardware efficiency and compact representation.

In 2008, St\u{a}nic\u{a} and Maitra formulated the following well-known conjecture:
\begin{quote}
\textbf{Conjecture 12 (\cite{stanica2008}):} \emph{There exist no homogeneous rotation-symmetric bent Boolean functions of algebraic degree $d > 2$ in any number of variables.}
\end{quote}

In degree $d = 2$, quadratic rotation-symmetric bent functions exist in all even dimensions; for instance, the antipodal function $f_0(x) = \sum_{i=0}^{n/2-1} x_i x_{i+n/2}$ is bent for every even $n$. For degrees $d > 2$, partial results have been achieved through various specialized techniques:
\begin{enumerate}
    \item \textbf{Monomial RSBFs:} Meng, Chen, and Fu~\cite{meng2010} proved that no monomial (single-orbit) homogeneous RSBF of degree $d > 2$ can be bent.
    \item \textbf{Cyclic distance constraints:} Meng et al.~\cite{meng2010} also showed non-existence when the maximum cyclic distance between consecutive indices of each monomial is at most $n/2$.
    \item \textbf{Prime-multiple dimensions $n = 2p$:} Cusick and Sanger~\cite{cusick2017} analyzed the Hadamard matrix $M = (\hat{f}(u \oplus v))$ and proved that for $n = 2p$ ($p$ an odd prime), every homogeneous RSBF that is bent must have even degree, thereby ruling out cubic HRSBFs for $n \in \{6, 10, 14, 22, 26, 34, \dots\}$.
    \item \textbf{Recent SANF and derivative bounds:} Sun, Shi, Liu, and Fu~\cite{sun2026} studied the short algebraic normal form (SANF), derivatives, and nonlinearity bounds, deriving partial non-existence results. Calder\'on-G\'omez et al.~\cite{calderon2024} established orbit-counting structural results.
\end{enumerate}

Despite these efforts, the cubic case ($d = 3$) has remained unresolved for composite factors, power-of-two dimensions ($n = 8, 16, 32, \dots$), and mixed families ($n = 4m, 8m, 16m$ for odd $m$). Table~\ref{tab:lit_comparison} contextualizes our contribution within the literature.

\begin{table}[htbp]
\centering
\small
\caption{Comparison of Non-existence Results for Cubic HRSBFs in the Literature.}
\label{tab:lit_comparison}
\begin{tabular}{llll}
\toprule
Reference & Method & Dimensions Covered & Scope / Hypotheses \\
\midrule
Meng et al. (2010)~\cite{meng2010} & Cyclic monomial distances & All $n$ & Monomials ($|\mathcal{O}|=1$), max distance $\le n/2$ \\
Zhang \& Gao (2013)~\cite{zhang2013} & 2-adic valuation \& SANF & Various & Specific conditions on monomial supports \\
Cusick \& Sanger (2017)~\cite{cusick2017} & Hadamard matrix entry $M_{0, 2^n-1}$ & $n = 2p$ ($p$ odd prime) & General cubic HRSBFs for $n=2p$ \\
Calder\'on-G\'omez et al. (2024)~\cite{calderon2024} & Short $k$-rotation orbits & Various & Orbit enumerations; conjecture open \\
Sun et al. (2026)~\cite{sun2026} & SANF / derivative bounds & Partial & Monomial support conditions \\
\textbf{This Work} & \textbf{Odd-order reduction + Ward divisibility} & $\mathbf{n \not\equiv 0 \pmod{32}}$ & \textbf{All $n \in \{2m, 4m, 8m, 16m\}$ ($\forall m$ odd)} \\
\bottomrule
\end{tabular}
\end{table}

\subsection{Our Contributions}
In this paper, we resolve the cubic St\u{a}nic\u{a}--Maitra conjecture for all even dimensions not divisible by 32. Our contributions are threefold:
\begin{enumerate}
    \item \textbf{Odd-Order Quadratic Reduction Lemma:} We prove an exact reduction lemma for linear automorphisms of odd order $m$ over $\F^n$: a $T$-invariant quadratic function $q$ is balanced on $\F^n$ if and only if its restriction to the fixed subspace $\mathrm{Fix}(T)$ is balanced. As a consequence, every directional derivative of a cubic bent function preserves its balance upon restriction, forcing the restriction of any cubic bent function to $\mathrm{Fix}(T)$ to remain bent.
    \item \textbf{Complete Orbit Filter and Structural Reduction:} For cyclic shifts $T = \sigma^t$ ($n = t \cdot m$, $t = 2^s$, $m$ odd), we show that all short orbits of length $l < t$ cancel modulo 2, and the antipodal quadratic orbit $q_{t/2}$ cancels identically. Since $f$ is homogeneous, no constant term is generated, and the restriction $g = f|_U$ belongs strictly to the linear space $\mathcal{F}_t = \mathrm{span}\{O^{(2)}_d, O^{(3)}_{a,b}, L\}$, preserving identical vanishing on the antipodal subspace $V_{t/2}$.
    \item \textbf{Spectral Divisibility Obstruction via Divisible Codes ($t = 16$):} We prove that the restricted complete-orbit space in 16 variables forms a 16-divisible binary linear code. Using an exact weight polarization expansion within the divisible codes framework of Ward~\cite{ward1981,ward2001}, we prove that every codeword satisfies $\wt(g) \equiv 0 \pmod{256}$, which algebraically forces $W_g(0) \equiv 0 \pmod{512} \ne \pm 256$, ruling out bentness for all $2^{43} \approx 8.79 \times 10^{12}$ functions (and $2^{42}$ in the subspace without $L$).
\end{enumerate}
Together with the analytical resolution of $t \in \{2, 4, 8\}$, this establishes the non-existence of cubic HRSBFs for all $n \in \{2m, 4m, 8m, 16m\}$ for all odd $m$. Notably, this covers $n = 2m$ where $m$ is composite (e.g., $n = 18, 30, 42, 50, \dots$), which was outside the reach of prime-based Hadamard arguments~\cite{cusick2017}, as well as the new families $n = 4m, 8m, 16m$.

%======================================================================
\section{Preliminaries}
\label{sec:prelim}
%======================================================================

Let $\V_n = \F^n$ denote the $n$-dimensional vector space over $\F$. A Boolean function $f$ in $n$ variables is a mapping $f: \V_n \to \F$. The Hamming weight of $f$, denoted $\wt(f)$, is the cardinality of its support $\{x \in \V_n : f(x) = 1\}$. A function $f$ is \emph{balanced} if $\wt(f) = 2^{n-1}$.

Every Boolean function has a unique representation as a multilinear polynomial over $\F$, its Algebraic Normal Form (ANF):
\begin{equation}
f(x) = \sum_{S \subseteq \{0, 1, \dots, n-1\}} c_S \prod_{i \in S} x_i, \quad c_S \in \F.
\end{equation}
The algebraic degree of $f$, denoted $\deg(f)$, is $\max\{|S| : c_S \ne 0\}$. The function $f$ is \emph{homogeneous of degree $d$} if $c_S = 0$ for all $|S| \ne d$.

The directional derivative of $f$ in direction $a \in \V_n$ is defined by:
\begin{equation}
D_a f(x) = f(x \oplus a) \oplus f(x).
\end{equation}
If $\deg(f) = d$, then $\deg(D_a f) \le d - 1$ for all $a \in \V_n$.

The Walsh transform of $f$ at $\omega \in \V_n$ is the discrete Fourier transform of its sign function $(-1)^{f(x)}$:
\begin{equation}
W_f(\omega) = \sum_{x \in \V_n} (-1)^{f(x) \oplus \omega \cdot x}.
\end{equation}
A function $f$ on $\V_n$ is \emph{bent} if $|W_f(\omega)| = 2^{n/2}$ for all $\omega \in \V_n$. Bent functions exist only when $n$ is even. By the well-known directional derivative characterization~\cite{meier1989,carlet2021}, $f$ is bent if and only if $D_a f$ is balanced on $\V_n$ for every non-zero direction $a \in \V_n \setminus \{0\}$.

A function $f$ is \emph{rotation-symmetric} (RSBF) if $f(\sigma(x)) = f(x)$ for all $x \in \V_n$, where $\sigma(x_0, x_1, \dots, x_{n-1}) = (x_1, \dots, x_{n-1}, x_0)$ is the cyclic shift operator. Under the action of the cyclic group $C_n = \langle \sigma \rangle$, the set of monomials of degree $d$ is partitioned into disjoint orbits. A homogeneous RSBF of degree $d$ is an $\F$-linear combination of the orbit sums:
\begin{equation}
f(x) = \sum_{j} c_j O_j(x), \quad O_j(x) = \sum_{m \in \mathcal{O}_j} m(x).
\end{equation}

%======================================================================
\section{The Odd-Order Quadratic Reduction Lemma}
\label{sec:reduction}
%======================================================================

Let $V = \V_n = \F^n$. Let $T: V \to V$ be an invertible linear automorphism of odd order $m$ ($\gcd(m, 2) = 1$). Let $U = \mathrm{Fix}(T) = \{x \in V : Tx = x\}$ denote the fixed subspace of $T$.

\begin{lemma}[Symplectic Decomposition and Exponential Sum Factorization]
\label{lem:quadratic_reduction}
Let $T \in \mathrm{GL}(V)$ have odd order $m$. Let $q: V \to \F$ be a quadratic Boolean function ($\deg(q) \le 2$) such that $q(0) = 0$ and $q(Tx) = q(x)$ for all $x \in V$. Then:
\begin{equation}
q \text{ is balanced on } V \iff q|_U \text{ is balanced on } U.
\end{equation}
\end{lemma}

\begin{proof}
Consider the polar form of $q$, defined by:
\begin{equation}
B(x, y) = q(x \oplus y) \oplus q(x) \oplus q(y).
\end{equation}
Since $\deg(q) \le 2$ and $q(0) = 0$, $B$ is a symmetric (alternating) bilinear form over $\F$. The $T$-invariance of $q$ implies that $B$ is $T$-invariant: $B(Tx, Ty) = B(x, y)$.

Define the projection operator $P: V \to V$ by:
\begin{equation}
P = \sum_{j=0}^{m-1} T^j.
\end{equation}
Since $m$ is odd, $m \equiv 1 \pmod 2$. Consequently, for any $x \in V$, $T(Px) = P(Tx) = Px$, so $\mathrm{Im}(P) \subseteq U$. Conversely, for $u \in U$, $Pu = \sum_{j=0}^{m-1} u = m \cdot u = u$. Thus, $P^2 = P$ and $\mathrm{Im}(P) = U$. 

This defines a direct sum decomposition $V = U \oplus K$, where $K = \ker(P)$. We now claim that $U$ and $K$ are orthogonal with respect to the bilinear form $B$. Indeed, let $u \in U$ and $k \in K$. Using the $T$-invariance of $B$ and the fact that $Tu = u$:
\begin{equation}
B(u, k) = \sum_{j=0}^{m-1} B(u, T^j k) = B\left(u, \sum_{j=0}^{m-1} T^j k\right) = B(u, Pk) = B(u, 0) = 0.
\end{equation}
Since $B(u, k) = 0$ for all $u \in U, k \in K$, we have $q(u \oplus k) = q(u) \oplus q(k) \oplus B(u, k) = q(u) \oplus q(k)$.

Therefore, the exponential sum $S_V(q) = \sum_{x \in V} (-1)^{q(x)}$ factors as:
\begin{equation}
S_V(q) = \left( \sum_{u \in U} (-1)^{q(u)} \right) \left( \sum_{k \in K} (-1)^{q(k)} \right) = S_U(q|_U) \cdot S_K(q|_K).
\end{equation}
It remains to show that $S_K(q|_K) \ne 0$. We compute the square of the exponential sum directly:
\begin{align}
S_K(q|_K)^2 &= \left( \sum_{x \in K} (-1)^{q(x)} \right) \left( \sum_{y \in K} (-1)^{q(y)} \right) = \sum_{x, y \in K} (-1)^{q(x) \oplus q(y)} \nonumber \\
&= \sum_{x, z \in K} (-1)^{q(x \oplus z) \oplus q(x)} = \sum_{z \in K} (-1)^{q(z)} \sum_{x \in K} (-1)^{B(x, z)},
\end{align}
where we substituted $y = x \oplus z$ and used $q(x \oplus z) \oplus q(x) = q(z) \oplus B(x, z)$. For each fixed $z \in K$, the mapping $x \mapsto B(x, z)$ is a linear form on $K$. By character orthogonality on $\F$, the inner sum satisfies:
\begin{equation}
\sum_{x \in K} (-1)^{B(x, z)} = \begin{cases} |K|, & \text{if } z \in \rad(B|_K), \\ 0, & \text{otherwise}. \end{cases}
\end{equation}
Therefore, letting $R_K = \rad(B|_K)$:
\begin{equation}
S_K(q|_K)^2 = |K| \sum_{z \in R_K} (-1)^{q(z)}.
\end{equation}
Since $B$ vanishes identically on $R_K$, the restriction $q|_{R_K}$ is a linear form. Using the $T$-invariance of $q$ and the fact that $m$ is odd:
\begin{equation}
q(z) = m \cdot q(z) = \sum_{j=0}^{m-1} q(T^j z) = q\left(\sum_{j=0}^{m-1} T^j z\right) = q(Pz).
\end{equation}
Since $z \in K = \ker(P)$, we have $Pz = 0$, and since $q(0) = 0$, we obtain $q(z) = 0$ for all $z \in R_K$. Consequently, $(-1)^{q(z)} = +1$ identically on $R_K$, yielding:
\begin{equation}
S_K(q|_K)^2 = |K| \cdot |R_K| > 0.
\end{equation}
Since $S_K(q|_K)^2 > 0$ over $\mathbb{R}$, we have $S_K(q|_K) \ne 0$. It follows immediately that $S_V(q) = 0 \iff S_U(q|_U) = 0$, which completes the proof.
\end{proof}

\begin{remark}[Algebraic Context and Literature Precedents]
The linear decomposition $V = U \oplus K$ under an odd-order automorphism group $G$ is rooted in classical representation theory over fields of characteristic 2: since $\gcd(|G|, 2) = 1$, the group algebra $\F[G]$ is semisimple by Maschke's theorem. Furthermore, the orthogonality $B(u, k) = 0$ is a standard property of invariant bilinear forms under group actions. Dickson's classical theorem on quadratic forms over $\F$ provides the foundational structure for the non-vanishing exponential sum $S_K(q|_K) \ne 0$, exploiting the fact that $q$ vanishes identically on the radical $R_K$ because $P(R_K) = 0$. 

However, the specific realization that $q$ is balanced on the full space if and only if its restriction to $\mathrm{Fix}(G)$ is balanced, and its application to transfer the bentness of cubic functions through the balance of their directional derivatives, is to the best of our knowledge new in the study of rotation-symmetric Boolean functions.
\end{remark}

\begin{remark}[Generalization to Arbitrary Odd-Order Groups]
The proof of Lemma~\ref{lem:quadratic_reduction} does not require $T$ to be cyclic. For any finite group $G \le \mathrm{GL}(V)$ of odd order $|G|$, the Reynolds operator $P = \sum_{g \in G} g$ satisfies $P^2 = P$ in characteristic 2 and yields $B(u, k) = 0$ for all $u \in \mathrm{Fix}(G), k \in \ker(P)$. Thus, the quadratic balance reduction holds for any odd-order linear group action.
\end{remark}

%======================================================================
\section{Restriction of Cubic Bent Functions}
\label{sec:cubic_bent}
%======================================================================

We now state the key consequence for bent Boolean functions of degree at most 3.

\begin{theorem}[Preservation of Bentness under Odd-Order Restriction]
\label{thm:bent_restriction}
Let $V = \F^n$ with $n$ even. Let $T \in \mathrm{GL}(V)$ have odd order $m$. Let $f: V \to \F$ be a Boolean function of degree $\deg(f) \le 3$ such that $f(Tx) = f(x)$ for all $x \in V$. If $f$ is bent, then its restriction $g = f|_U$ to the fixed subspace $U = \mathrm{Fix}(T)$ is bent on $U$, provided $\dim(U)$ is even.
\end{theorem}

\begin{proof}
Let $a \in U \setminus \{0\}$. Consider the directional derivative:
\begin{equation}
D_a f(x) = f(x \oplus a) \oplus f(x).
\end{equation}
Since $\deg(f) \le 3$, we have $\deg(D_a f) \le 2$. Furthermore, since $Ta = a$ and $f \circ T = f$:
\begin{equation}
D_a f(Tx) = f(Tx \oplus a) \oplus f(Tx) = f(T(x \oplus a)) \oplus f(Tx) = f(x \oplus a) \oplus f(x) = D_a f(x).
\end{equation}
Thus, $D_a f$ is $T$-invariant.

To apply Lemma~\ref{lem:quadratic_reduction}, we normalize the constant term by setting:
\begin{equation}
q(x) = D_a f(x) \oplus D_a f(0).
\end{equation}
Then $q(0) = 0$, $\deg(q) \le 2$, and $q$ is $T$-invariant. Since $f$ is bent, $D_a f$ is balanced on $V$. Adding the constant $D_a f(0)$ merely swaps the preimages $0 \leftrightarrow 1$ if $D_a f(0) = 1$, so $q$ is balanced on $V$.

By Lemma~\ref{lem:quadratic_reduction}, $q|_U$ is balanced on $U$. Consequently, $(D_a f)|_U = q|_U \oplus D_a f(0)$ is balanced on $U$. But for all $u \in U$:
\begin{equation}
(D_a f)|_U(u) = f(u \oplus a) \oplus f(u) = D_a(f|_U)(u) = D_a g(u).
\end{equation}
Therefore, every non-zero directional derivative $D_a g$ for $a \in U \setminus \{0\}$ is balanced on $U$. By the directional derivative characterization~\cite{meier1989,carlet2021}, $g$ is bent on $U$.
\end{proof}

%======================================================================
\section{Complete Orbit Filtering and Affine Invariance}
\label{sec:orbit_filter}
%======================================================================

Let $n = t \cdot m$, where $t = 2^s$ ($s \ge 1$) and $m$ is odd. Let $T = \sigma^t$ be the cyclic shift by $t$ positions. The order of $T$ is $m$. The fixed subspace $U = \mathrm{Fix}(T)$ consists of all $x \in \V_n$ formed by $m$ repetitions of a binary block of length $t$:
\begin{equation}
x = (y, y, \dots, y) \in \V_n, \quad y = (y_0, y_1, \dots, y_{t-1}) \in \V_t.
\end{equation}
Thus, $U \cong \V_t$ has dimension $t = 2^s$, which is even for $s \ge 1$.

\subsection{Parity Cancellation of Short Orbits}
Let $f$ be a homogeneous cubic RSBF in $n$ variables:
\begin{equation}
f(x) = \sum_{j} c_j O_j(x), \quad O_j(x) = \sum_{k=0}^{L-1} \sigma^k(x_0 x_a x_b).
\end{equation}
The length $L$ of the orbit under $\sigma$ divides $n$. Under the cyclic action on triplets of coordinates, the stabilizer of $\{0, a, b\}$ has cardinality $h \in \{1, 3\}$. Hence $L \in \{n, n/3\}$. In both cases, the largest power of 2 dividing $L$ is $t = 2^s$.

\begin{proposition}[Complete Orbit Filter]
\label{prop:complete_orbit_filter}
Upon restriction to $U = \mathrm{Fix}(\sigma^t) \cong \V_t$, only orbits of full length $t$ survive in $\F$.
\end{proposition}

\begin{proof}
Let $\mu(x) = x_i x_j x_k$ be a cubic monomial in $V$. The subgroup $H \le \mathbb{Z}_n$ of cyclic shifts fixing the index set $\{i, j, k\}$ acts freely on $\{i, j, k\}$, because any non-zero cyclic translation on $\mathbb{Z}_n$ has no fixed points. Therefore, every orbit of the action of $H$ on $\{i, j, k\}$ has cardinality $|H|$, which implies that $|H|$ divides $|\{i, j, k\}| = 3$. Since 3 is prime, $|H| \in \{1, 3\}$.

Since $n = t \cdot m$ with $m$ odd and $|H| \in \{1, 3\}$ is odd, the ratio $m' = m / |H|$ is an odd integer:
\begin{equation}
\frac{L}{t} = \frac{n / |H|}{t} = \frac{m}{|H|} \equiv 1 \pmod 2.
\end{equation}
Decomposing the cyclic summation of length $L = t \cdot m'$ into $m'$ blocks of size $t$:
\begin{equation}
\sum_{k=0}^{L-1} \sigma^k(\mu)(x) = \sum_{r=0}^{m'-1} \sum_{j=0}^{t-1} \sigma^{j + r t}(\mu)(x).
\end{equation}
Upon restriction to $U = \mathrm{Fix}(\sigma^t)$, the shift $\sigma^{rt}$ acts trivially. Since $m'$ is an odd integer, the summation in $\F$ reduces to:
\begin{equation}
\left( \sum_{k=0}^{L-1} \sigma^k(\mu) \right) \Bigg|_U = m' \sum_{j=0}^{t-1} \sigma^j(\mu|_U) = \sum_{j=0}^{t-1} \sigma^j(\mu|_U).
\end{equation}
That is, the restriction of the original orbit collapses to the cyclic shift sum over $\mathbb{Z}_t$ of the reduced monomial $\mu|_U$.

Now, under the cyclic shift of $t$ coordinates, the orbit of $\mu|_U$ has length $l$, where $l \mid t = 2^s$. The cyclic sum of $t$ terms covers this reduced orbit exactly $t/l$ times. If $l < t$, then $l = 2^r$ with $r < s$, so $t/l = 2^{s-r} \ge 2 \equiv 0 \pmod 2$.
Therefore, in $\F$, every short orbit with $l < t$ cancels out completely. Only orbits of full length $l = t$ survive.
\end{proof}

\begin{lemma}[Exact Parity Cancellation of the Antipodal Orbit $q_{t/2}$]
\label{lem:cancellation_qt2}
Let $n = t \cdot m$ with $t = 2^s$ and $m$ odd. The antipodal quadratic orbit $q_{t/2}(y) = \sum_{r=0}^{t/2-1} y_r y_{r+t/2}$ cannot appear in the restriction $g = f|_U$ of any homogeneous cubic RSBF $f$: its coefficient in $\F$ is strictly zero.
\end{lemma}

\begin{proof}
A cubic monomial $\mu(x) = x_i x_j x_k$ in $\V_n$ restricts to a quadratic term on $U = \mathrm{Fix}(\sigma^t)$ if and only if two of its indices are congruent modulo $t$ and the third is distinct: $i \equiv j \pmod t$ with $k \not\equiv i \pmod t$. Let $a = i \bmod t$ and $b = k \bmod t$. Then $\mu|_U = y_a y_b$. For this term to belong to the antipodal quadratic orbit, the two indices must differ by $t/2$ modulo $t$, so $b \equiv a + t/2 \pmod t$.

Under the cyclic shifts $\sigma^s$ of $\V_n$, the monomial $\sigma^s(\mu)$ restricts to $y_{(a+s)\bmod t} y_{(a+s+t/2)\bmod t}$. Consider a fixed antipodal pair $\{r, r + t/2\}$ of $q_{t/2}$. Because $\{r + t/2, (r + t/2) + t/2\} \equiv \{r, r + t/2\} \pmod t$, this pair is invariant under translation by $t/2$ in $\mathbb{Z}_t$. Therefore, a shift $\sigma^s(\mu)$ restricts to the pair $\{r, r + t/2\}$ if and only if $s \equiv r - a \pmod{t/2}$.

The orbit of $\mu$ has length $L = n / |\mathrm{Stab}(\mu)| = t \cdot m / |\mathrm{Stab}(\mu)|$. Because $L$ is a multiple of $t/2$, the number of shifts $s \in \{0, \dots, L-1\}$ satisfying the congruence $s \equiv r - a \pmod{t/2}$ is exactly:
\begin{equation}
\frac{L}{t/2} = \frac{t \cdot m / |\mathrm{Stab}(\mu)|}{t/2} = \frac{2m}{|\mathrm{Stab}(\mu)|}.
\end{equation}
The cyclic stabilizer $H = \mathrm{Stab}(\mu) \le \mathbb{Z}_n$ acts freely on the index set $\{i, j, k\}$ (since non-zero translations in $\mathbb{Z}_n$ have no fixed points). Hence $|H|$ divides $|\{i, j, k\}| = 3$. Since 3 and $m$ are odd, $|H| \in \{1, 3\}$ is an odd divisor of $m$, so $m / |H|$ is an integer. Thus:
\begin{equation}
\frac{2m}{|H|} = 2 \left( \frac{m}{|H|} \right) \equiv 0 \pmod 2.
\end{equation}
In $\F$, this multiplicity cancels completely ($2 \equiv 0$). Summing over all cubic orbits in $f$, the coefficient of each pair $y_r y_{r+t/2}$ is $0 \pmod 2$. Hence, $q_{t/2}$ never appears in $g = f|_U$.
\end{proof}

\begin{proposition}[Appearance of the All-Ones Linear Form $L$]
\label{prop:linear_form_L}
Let $n = t \cdot m$ with $m \ge 3$ odd. Any cubic monomial $\mu(x) = x_i x_j x_k$ with $i \equiv j \equiv k \pmod t$ restricts on $U$ to $y_a^3 = y_a$, where $a = i \bmod t$. The number of shifts in the orbit of $\mu$ that project to a fixed residue $r \in \mathbb{Z}_t$ is:
\begin{equation}
\frac{L}{t} = \frac{m}{|\mathrm{Stab}(\mu)|} \equiv 1 \pmod 2,
\end{equation}
since both $m$ and $|\mathrm{Stab}(\mu)| \in \{1, 3\}$ are odd. Consequently, each such orbit contributes with coefficient 1 to every coordinate $y_r$, producing the full all-ones linear form:
\begin{equation}
L(y) = \sum_{r=0}^{t-1} y_r.
\end{equation}
All lower-degree terms generated upon restriction are therefore spanned strictly by the all-ones linear form $L(y)$.
\end{proposition}

\subsection{Absence of Constant Terms and the Restricted Space $\mathcal{F}_t$}
Since the original function $f$ is strictly homogeneous of degree 3, we have $f(0) = 0$. Under the coordinate restriction mapping $x_i \mapsto y_{i \bmod t}$, every cubic monomial $x_i x_j x_k$ yields either a cubic term (if the three indices are distinct modulo $t$), a quadratic term (if two indices collide), or a linear term (if all three indices collide). In particular, no degree-0 (constant) terms can ever be generated upon restriction.

Consequently, the restricted function satisfies $g(0) = f(0) = 0$, containing no constant term ($c_0 = 0$). The function $g = f|_U$ belongs strictly to the linear space:
\begin{equation}
\mathcal{F}_t = \mathrm{span}\{O^{(2)}_d, O^{(3)}_{a,b}, L\}.
\end{equation}
The strict absence of a constant term is mathematically essential: it ensures that the antipodal vanishing property $g|_{V_{t/2}} \equiv 0$ (established in Section~\ref{sec:t16}) holds without exception, as any non-zero constant $c_0 = 1$ would directly violate vanishing on $V_{t/2}$.

%======================================================================
\section{Exclusion of the Base Cases: $t \in \{2, 4, 8\}$}
\label{sec:base_cases}
%======================================================================

We now examine the reduced spaces $\mathcal{F}_t$ for $t \in \{2, 4, 8\}$ including the linear form $L$.

\begin{proposition}[Case $t = 2$ ($n = 2m$)]
\label{prop:t2}
In $t = 2$ variables, there are no complete orbits of degree 2 or 3. The space $\mathcal{F}_2$ is spanned solely by $L(y) = y_0 \oplus y_1$, containing $2^1 = 2$ functions ($g \equiv 0$ and $g = L$). Their Walsh values at zero are $W_g(0) \in \{0, 4\}$. Since bentness requires $|W_g(0)| = 2^{2/2} = 2$, no function in $\mathcal{F}_2$ is bent.
\end{proposition}

\begin{proposition}[Case $t = 4$ ($n = 4m$)]
\label{prop:t4}
In $t = 4$ variables, the generator set comprises 1 complete quadratic orbit $O_1^{(2)}$, 1 complete cubic orbit $O_1^{(3)}$, and the linear form $L(y) = \sum_{i=0}^3 y_i$. The space $\mathcal{F}_4$ contains $2^3 = 8$ functions. Exhaustive computation yields:
\begin{equation}
W_g(0) \in \{-8, 0, 8, 16\}.
\end{equation}
Since bentness in 4 variables requires $|W_g(0)| = 2^{4/2} = 4$, none of these 8 functions is bent.
\end{proposition}

\begin{proposition}[Case $t = 8$ ($n = 8m$)]
\label{prop:t8}
In $t = 8$ variables, there are 3 complete quadratic orbits ($d \in \{1, 2, 3\}$), 7 complete cubic orbits, and the linear form $L(y) = \sum_{i=0}^7 y_i$, spanning a space $\mathcal{F}_8$ of dimension 11 ($2^{11} = 2048$ functions). For all $g \in \mathcal{F}_8$, exact computation yields:
\begin{equation}
W_g(0) \in \{-192, -160, -128, -96, -64, -32, 0, 32, 64, 96, 128, 160, 192, 256\} \subset 32\Z.
\end{equation}
Since bentness in 8 variables requires $|W_g(0)| = 2^{8/2} = 16$, and $\pm 16 \not\equiv 0 \pmod{32}$, no function in $\mathcal{F}_8$ is bent.
\end{proposition}

%======================================================================
\section{The 16-Divisibility Obstruction via Ward's Theorem ($t = 16$)}
\label{sec:t16}
%======================================================================

In $t = 16$ variables, the complete-orbit space $\mathcal{F}_{16}$ is spanned by:
\begin{itemize}
    \item 7 complete quadratic orbits: $O_d(y) = \sum_{i=0}^{15} y_i y_{i+d}$ for $d \in \{1, 2, 3, 4, 5, 6, 7\}$.
    \item 35 complete cubic orbits: $O_{a,b}(y) = \sum_{i=0}^{15} y_i y_{i+a} y_{i+b}$ for valid cyclic gap-partition representatives $1 \le a < b < 16$.
    \item 1 all-ones linear form: $L(y) = \sum_{i=0}^{15} y_i$.
\end{itemize}
Thus, $\dim(\mathcal{F}_{16}) = 7 + 35 + 1 = 43$, representing $2^{43} \approx 8.79 \times 10^{12}$ candidate functions. Brute-force spectral evaluation is computationally intractable; we resolve this case analytically via coding theory.

\subsection{The Antipodal Subspace Vanishing Theorem}
Let $V_8 = \{y \in \V_{16} : y_{i+8} = y_i\}$ denote the antipodal subspace of dimension 8.

\begin{lemma}[Antipodal Vanishing]
\label{lem:antipodal_vanishing}
Every function $g \in \mathcal{F}_{16}$ satisfies $g(y) = 0$ identically for all $y \in V_8$.
\end{lemma}

\begin{proof}
Every complete orbit in $\mathcal{F}_{16}$ has length 16 and can be written as $\sum_{i=0}^7 (m_i(y) \oplus m_{i+8}(y))$. For $y \in V_8$, $y_{j+8} = y_j$ for all $j$, so $m_{i+8}(y) = m_i(y)$ identically. Thus, each pair cancels modulo 2, and $g|_{V_8} \equiv 0$.
\end{proof}

\subsection{Code Formulation and Weight Polarization}
Since $g|_{V_8} \equiv 0$, the non-zero values of $g$ lie entirely in $\V_{16} \setminus V_8$, which consists of $2^{16} - 2^8 = 65280$ vectors. Under the cyclic shift $C_{16}$, these vectors partition into exactly $65280 / 16 = 4080$ orbits of length 16.

Let $\mathcal{O}_1, \dots, \mathcal{O}_{4080}$ denote these orbits, with representatives $r_1, \dots, r_{4080}$. For each basis orbit $v_j \in \mathcal{F}_{16}$ ($1 \le j \le 42$), define its indicator codeword $w_j \in \F^{4080}$ by:
\begin{equation}
(w_j)_k = v_j(r_k), \quad 1 \le k \le 4080.
\end{equation}
Let $\mathcal{C} \subset \F^{4080}$ be the binary linear code generated by $\{w_1, \dots, w_{42}\}$. For any function $g = \sum_{j \in I} v_j$, its indicator codeword is $c = \sum_{j \in I} w_j \in \mathcal{C}$, and its Hamming weight on $\V_{16}$ is:
\begin{equation}
\wt(g) = 16 \cdot \wt(c).
\end{equation}

To determine the divisibility of $\wt(c)$, we employ the weight polarization expansion for binary codes within the divisible codes framework of Ward~\cite{ward1981,ward2001}:

\begin{theorem}[Weight Polarization Expansion Modulo 16]
\label{thm:ward_expansion}
Let $\mathcal{C}$ be a binary linear code spanned by generators $w_1, \dots, w_k$. For any subset $I \subseteq \{1, \dots, k\}$, the Hamming weight of $c = \sum_{i \in I} w_i$ satisfies:
\begin{equation}
\label{eq:ward_formula}
\wt(c) = \sum_{J \subseteq I, J \ne \emptyset} (-2)^{|J|-1} \wt\left( \bigcap_{j \in J} w_j \right).
\end{equation}
In particular, modulo 16:
\begin{equation}
\label{eq:ward_mod16}
\wt(c) \equiv \sum_{i} \wt(w_i) - 2 \sum_{i < j} \wt(w_i \cap w_j) + 4 \sum_{i < j < l} \wt(w_i \cap w_j \cap w_l) - 8 \sum_{i < j < l < p} \wt(w_i \cap w_j \cap w_l \cap w_p) \pmod{16},
\end{equation}
since $(-2)^{|J|-1} \equiv 0 \pmod{16}$ for all $|J| \ge 5$.
\end{theorem}

\begin{proof}
For any coordinate $x$, let $m(x) = |\{i \in I : (w_i)_x = 1\}|$ be the number of active generators at $x$. The coordinate $x$ belongs to $\supp(c)$ if and only if $m(x)$ is odd. Its contribution to the right-hand side of~\eqref{eq:ward_formula} is:
\begin{equation}
\sum_{k=1}^{m(x)} \binom{m(x)}{k} (-2)^{k-1} = -\frac{1}{2} \left[ (1 - 2)^{m(x)} - 1 \right] = \frac{1 - (-1)^{m(x)}}{2} = \begin{cases} 1, & m(x) \text{ odd}, \\ 0, & m(x) \text{ even}. \end{cases}
\end{equation}
Summing over all coordinates establishes~\eqref{eq:ward_formula}. Reducing modulo 16 eliminates all terms with $|J| \ge 5$ because $(-2)^4 = 16 \equiv 0 \pmod{16}$.
\end{proof}

\begin{theorem}[Spectral Obstruction for $t = 16$]
\label{thm:t16_obstruction}
For every function $g \in \mathcal{F}_{16}$, $\wt(g) \equiv 0 \pmod{256}$, and consequently $W_g(0) \equiv 0 \pmod{512}$. Therefore, no function in $\mathcal{F}_{16}$ is bent.
\end{theorem}

\begin{proof}
By Theorem~\ref{thm:ward_expansion}, $\wt(c) \equiv 0 \pmod{16}$ holds for all $c \in \mathcal{C}$ if the following four conditions are satisfied:
\begin{enumerate}
    \item $\wt(w_i) \equiv 0 \pmod{16}$ for all 43 generators (35 cubic, 7 quadratic, and $L$).
    \item $\wt(w_i \cap w_j) \equiv 0 \pmod 8$ for all $\binom{43}{2} = 903$ generator pairs.
    \item $\wt(w_i \cap w_j \cap w_l) \equiv 0 \pmod 4$ for all $\binom{43}{3} = 12341$ generator triples.
    \item $\wt(w_i \cap w_j \cap w_l \cap w_p) \equiv 0 \pmod 2$ for all $\binom{43}{4} = 123410$ generator quadruples.
\end{enumerate}
We performed an exact computation of all $136697$ intersection weights on the indicator matrix of $\mathcal{C}$ in $\F^{4080}$ (as well as the $124313$ subsets of the 42-generator subspace without $L$). Every condition is satisfied with zero violations (see Table~\ref{tab:ward_audit}).

\begin{table}[htbp]
\centering
\caption{Audit of Ward divisibility conditions for the full complete-orbit code $\mathcal{C} \subset \F^{4080}$ (43 generators).}
\label{tab:ward_audit}
\begin{tabular}{lccc}
\toprule
Intersection Order & Number of Subsets & Required Congruence & Violations \\
\midrule
Singletons ($k = 1$) & $43$ & $\wt(w_i) \equiv 0 \pmod{16}$ & $0$ \\
Pairs ($k = 2$) & $903$ & $\wt(w_i \cap w_j) \equiv 0 \pmod 8$ & $0$ \\
Triples ($k = 3$) & $12341$ & $\wt(w_i \cap w_j \cap w_l) \equiv 0 \pmod 4$ & $0$ \\
Quadruples ($k = 4$) & $123410$ & $\wt(w_i \cap w_j \cap w_l \cap w_p) \equiv 0 \pmod 2$ & $0$ \\
\midrule
\textbf{Total Subsets} & $\mathbf{136697}$ & (Subcase without $L$: $124313$) & $\mathbf{0}$ \\
\bottomrule
\end{tabular}
\end{table}

Thus, $\wt(c) \equiv 0 \pmod{16}$ for every codeword $c \in \mathcal{C}$. It follows that:
\begin{equation}
\wt(g) = 16 \cdot \wt(c) \equiv 0 \pmod{256}.
\end{equation}
The Walsh transform at $\omega = 0$ is:
\begin{equation}
W_g(0) = 2^{16} - 2\,\wt(g) = 65536 - 2(256 k) \equiv 0 \pmod{512}.
\end{equation}
For $g$ to be bent on $\V_{16}$, it must satisfy $|W_g(0)| = 2^{16/2} = 2^8 = 256$. Since $\pm 256 \not\equiv 0 \pmod{512}$, no function in $\mathcal{F}_{16}$ can be bent.
\end{proof}

\subsection{Theoretical Structure: Circulant Rank Deficit and Coset Compression}
To provide deeper structural insight into why the 16-divisibility conditions hold for the generators and their intersections, we establish two properties of the complete orbits.

\begin{theorem}[Exact Weight and Spectral Value of Quadratic Orbits via Transfer Matrix]
\label{thm:exact_quadratic_weight}
Let $t = 2^s$ with $s \ge 2$, and let $O_d(y) = \sum_{i=0}^{t-1} y_i y_{i+d}$ be a complete quadratic orbit with $1 \le d < t/2$. Let $a = v_2(d)$. Then the Walsh coefficient at zero and the Hamming weight of $O_d$ are given exactly by:
\begin{equation}
W_{O_d}(0) = + 2^{t/2 + 2^a}, \qquad \wt(O_d) = 2^{t-1} - 2^{t/2 + 2^a - 1}.
\end{equation}
In particular, the 2-adic valuation of the Hamming weight satisfies:
\begin{equation}
v_2(\wt(O_d)) = \frac{t}{2} + 2^a - 1 \ge \frac{t}{2}.
\end{equation}
\end{theorem}

\begin{proof}
The coordinate shift graph on indices $\{0, 1, \dots, t-1\}$ defined by edges $(i, (i+d) \bmod t)$ partitions into $c = \gcd(t, d) = 2^a$ disjoint cycles of length $L = t/c = 2^{s-a}$. Since $1 \le d < t/2$, we have $L \ge 4$ and $L$ is even.

Along any single cycle of length $L$, the exponential sum of the periodic 2-chain is given by the trace of the $L$-th power of the $2 \times 2$ transfer matrix:
\begin{equation}
\sum_{z \in \F^L} (-1)^{\sum_{i=0}^{L-1} z_i z_{i+1}} = \operatorname{tr}(A^L), \quad \text{where } A = \begin{pmatrix} 1 & 1 \\ 1 & -1 \end{pmatrix}.
\end{equation}
Since $A^2 = 2 I_2$ and $L$ is even, $A^L = (A^2)^{L/2} = (2 I_2)^{L/2} = 2^{L/2} I_2$. Therefore, the trace along a single cycle is:
\begin{equation}
\operatorname{tr}(A^L) = \operatorname{tr}(2^{L/2} I_2) = 2 \cdot 2^{L/2} = 2^{L/2 + 1}.
\end{equation}
Because the $c = 2^a$ cycles involve disjoint sets of variables, the complete exponential sum factors over the cycles:
\begin{equation}
W_{O_d}(0) = \prod_{j=1}^c \operatorname{tr}(A^L) = \left( 2^{L/2 + 1} \right)^c = 2^{c(L/2 + 1)} = 2^{t/2 + c} = 2^{t/2 + 2^a}.
\end{equation}
This explicitly fixes the positive sign. The Hamming weight is then:
\begin{equation}
\wt(O_d) = \frac{2^t - W_{O_d}(0)}{2} = 2^{t-1} - 2^{t/2 + 2^a - 1} = 2^{t/2 + 2^a - 1} \left( 2^{t/2 - 2^a} - 1 \right).
\end{equation}
Since $t/2 - 2^a \ge 1$, the factor $2^{t/2 - 2^a} - 1$ is odd, establishing that $v_2(\wt(O_d)) = t/2 + 2^a - 1 \ge t/2$.
\end{proof}

\begin{theorem}[Coset Degree Compression]
\label{thm:coset_compression}
Let $g \in \mathcal{F}_t$ with $t = 2^s$. For any coset transversal vector $w \in \F^{t/2}$, define the coset restriction $g_w: \F^{t/2} \to \F$ by $g_w(y) = g(y \oplus w, y)$. Then $\deg(g_w) \le 2$.
\end{theorem}

\begin{proof}
Let $O_{a,b}(x) = \sum_{i=0}^{t-1} x_i x_{i+a} x_{i+b}$ be any complete cubic orbit. Under the affine substitution $x = (y \oplus w, y)$ with a fixed parameter $w \in \F^{t/2}$, the ANF coefficient of any cubic monomial $y_i y_j y_k$ (for distinct indices) is identically equal to its coefficient in the unshifted restriction $x = (y, y) \in V_{t/2}$, because the constant vector $w$ can only contribute to terms of lower degree in $y$. By Lemma~\ref{lem:antipodal_vanishing}, the restriction $O_{a,b}|_{V_{t/2}}$ vanishes identically, so all degree-3 coefficients vanish. Furthermore, coordinate identifications $y_i^2 = y_i$ resulting from index collisions modulo $t/2$ can only reduce monomial degree. Consequently, no degree-3 terms survive, establishing that $\deg(g_w) \le 2$ for all $w \in \F^{t/2}$. Structurally, this ensures that the restriction of $g$ to each affine coset of the antipodal subspace behaves as a quadratic form, providing algebraic motivation for the spectral divisibility.
\end{proof}

\subsection{Computational Protocol and Reproducibility}
To ensure complete scientific transparency and independent reproducibility:
\begin{enumerate}
    \item \textbf{Orbit Enumeration:} Orbits under $C_{16}$ are enumerated by minimal lexicographic index representatives. Out of $\binom{16}{2} = 120$ quadratic pairs, exactly 7 form full orbits of length 16 (the antipodal distance $d = 8$ forms an orbit of length 8). Out of $\binom{16}{3} = 560$ cubic triplets, exactly 35 form full orbits of length 16 (none have shorter periods).
    \item \textbf{Matrix Construction:} The $42$ generator orbits are evaluated over the $4080$ orbit representatives of $\V_{16} \setminus V_8$, generating a binary matrix $G \in \F^{42 \times 4080}$.
    \item \textbf{Exact Subset Verification:} Rather than computing weights of all $2^{42}$ codewords, Theorem~\ref{thm:ward_expansion} establishes that checking the 124,313 intersection weights is mathematically equivalent to checking all $2^{42}$ codewords.
    \item \textbf{Execution:} The verification protocol is implemented in a standalone, self-contained Python script (\texttt{verificador\_independente.py}) requiring zero external libraries. On a standard CPU, the complete suite executes in 1.1 seconds.
\end{enumerate}

%======================================================================
\section{The Main Theorem}
\label{sec:main_theorem}
%======================================================================

We now assemble the complete result.

\begin{theorem}[Main Theorem]
\label{thm:main}
Let $n$ be an even positive integer such that $n \not\equiv 0 \pmod{32}$. Then there exists no cubic homogeneous rotation-symmetric bent Boolean function in $n$ variables.
\end{theorem}

\begin{proof}
Let $n = 2^s \cdot m$, where $m$ is odd. Since $n$ is even and $n \not\equiv 0 \pmod{32}$, the 2-adic valuation satisfies $v_2(n) = s \in \{1, 2, 3, 4\}$. Set $t = 2^s \in \{2, 4, 8, 16\}$.

Suppose there exists a cubic homogeneous RSBF $f$ in $n$ variables that is bent. Let $T = \sigma^t$. The order of $T$ is $m$, which is odd. By Theorem~\ref{thm:bent_restriction}, the restricted function $g = f|_U$ to the fixed subspace $U = \mathrm{Fix}(\sigma^t) \cong \V_t$ is bent on $\V_t$.

By Proposition~\ref{prop:complete_orbit_filter}, all short orbits of length $l < t$ cancel modulo 2, and by Lemma~\ref{lem:cancellation_qt2}, the antipodal quadratic orbit $q_{t/2}$ has coefficient 0. Furthermore, because $f$ is homogeneous of degree 3, no constant term can arise upon restriction, so the restricted function $g = f|_U$ belongs strictly to the linear space $\mathcal{F}_t = \mathrm{span}\{O^{(2)}_d, O^{(3)}_{a,b}, L\}$.

However:
\begin{itemize}
    \item If $s = 1$ ($t = 2$), Proposition~\ref{prop:t2} shows that $\mathcal{F}_2$ contains no bent functions.
    \item If $s = 2$ ($t = 4$), Proposition~\ref{prop:t4} shows that $\mathcal{F}_4$ contains no bent functions.
    \item If $s = 3$ ($t = 8$), Proposition~\ref{prop:t8} shows that $\mathcal{F}_8$ contains no bent functions.
    \item If $s = 4$ ($t = 16$), Theorem~\ref{thm:t16_obstruction} shows that $\mathcal{F}_{16}$ contains no bent functions.
\end{itemize}
This contradiction establishes that no such bent function $f$ exists.
\end{proof}

%======================================================================
\section{Exhaustive Computational Exclusions in Low Dimensions}
\label{sec:low_dim}
%======================================================================

To complement our analytical reduction theorems and provide independent verification, we performed exhaustive end-to-end searches across all low-dimensional candidate spaces.

\begin{proposition}[Low-Dimensional Cubic Homogeneous RSBFs ($n \le 12$)]
\label{prop:cubic_low_dim}
There exist no cubic homogeneous rotation-symmetric bent Boolean functions in $n \in \{4, 6, 8, 10, 12\}$ variables:
\begin{itemize}
    \item $n = 4$: 1 orbit, 2 functions $\to$ 0 bent.
    \item $n = 6$: 4 orbits, $2^4 = 16$ functions evaluated via 64-point FWHT $\to$ 0 bent ($< 0.01$ s).
    \item $n = 8$: 7 orbits, $2^7 = 128$ functions evaluated via 256-point FWHT $\to$ 0 bent ($< 0.05$ s).
    \item $n = 10$: 12 orbits, $2^{12} = 4096$ functions evaluated via 1024-point FWHT $\to$ 0 bent ($< 0.3$ s).
    \item $n = 12$: 19 orbits, spanning $2^{19} = 524{,}288$ functions. Testing the necessary zero-frequency condition $|W_f(0)| = 2^{12/2} = 64$ ($\wt(f) \in \{2016, 2080\}$) selects exactly $32{,}878$ candidates; verifying directional derivative balance and full 4096-point FWHT confirms that exactly 0 functions are bent ($13.1$ s).
\end{itemize}
\end{proposition}

\begin{theorem}[Exclusion of Quartic HRSBFs for $n = 8$ and $n = 10$]
\label{thm:quartic}
There exist no quartic homogeneous rotation-symmetric bent Boolean functions in 8 or 10 variables:
\begin{enumerate}
    \item \textbf{Case $n = 8, d = 4$:} There are 10 rotation-symmetric orbits of degree 4, generating $2^{10} = 1024$ candidate functions. Exhaustive evaluation of the full 256-point fast Walsh-Hadamard transform (FWHT) across all 1024 functions confirms that exactly 0 functions are bent (execution time $< 0.3$ seconds).
    \item \textbf{Case $n = 10, d = 4$:} There are 22 rotation-symmetric orbits of degree 4, spanning $2^{22} = 4{,}194{,}304$ candidate functions. Under the cyclic group $C_{10}$, the space $\V_{10}$ partitions into 108 orbits. Using an incremental Gray-code traversal over the 108-dimensional orbit basis:
    \begin{itemize}
        \item Testing the necessary spectral condition $|W_f(0)| = 2^{10/2} = 32$ ($\wt(f) \in \{496, 528\}$) filters the space down to exactly $78{,}116$ candidate functions.
        \item Testing the balance of the directional derivatives $D_e f$ along the standard unit directions $e \in \{e_0, \dots, e_9\}$ further prunes the candidates to exactly $1{,}706$ surviving functions.
        \item Computing the full 1024-point Walsh-Hadamard transform on all $1{,}706$ survivors reveals that exactly 0 functions are bent (total execution time $< 50$ seconds).
    \end{itemize}
\end{enumerate}
\end{theorem}
All verification routines are implemented in \texttt{verificador\_independente.py} and execute deterministically in seconds.

%======================================================================
\section{Broader Impacts and Cross-Disciplinary Applications}
\label{sec:applications}
%======================================================================

Beyond the resolution of the St\u{a}nic\u{a}--Maitra conjecture for $n \not\equiv 0 \pmod{32}$, the mathematical framework developed in this work---uniting odd-order symplectic reductions, projective orbit filtering, and Ward's divisible codes---has immediate ramifications across several allied disciplines:

\subsection{Cryptographic S-Box Design and Stream Cipher Security}
In symmetric-key cryptography, bent functions, almost bent (AB) functions, and almost perfect nonlinear (APN) functions serve as foundational components for substitution boxes (S-boxes) and nonlinear filter generators in stream ciphers (e.g., Grain, Trivium, and SNOW) and lightweight block ciphers (e.g., LowMC, ASCON, and Clyde). A persistent design objective has been to utilize rotation-symmetric functions due to their minimal hardware implementation footprint and pipeline efficiency.
Our Main Theorem proves definitively that designers cannot exploit cubic homogeneous rotation-symmetric functions as bent components, thereby establishing sharp architectural bounds for cryptographic primitives. Furthermore, the Odd-Order Quadratic Reduction Lemma (Lemma~\ref{lem:quadratic_reduction}) provides a universal tool to analyze the differential uniformity and algebraic immunity of vectorial RSBFs under sub-register decimation, allowing cryptanalysts to evaluate security against higher-order differential attacks without full-scale exponential sum computation.

\subsection{Quantum Error-Correcting Codes and Transversal Quantum Gates}
In fault-tolerant quantum computing, the implementation of transversal non-Clifford gates (specifically the $\pi/8$ phase gate $T = \mathrm{diag}(1, e^{i\pi/4})$) is essential for achieving universal quantum computation without prohibitive magic-state distillation overhead. By the Eastin-Knill theorem, transversal gates cannot achieve universality; however, Calderbank-Shor-Steane (CSS) quantum codes constructed from doubly-even, triply-even, or higher $2^k$-divisible binary linear codes enable transversal implementations of high-level Clifford hierarchy gates.
Our construction in Theorem~\ref{thm:t16_obstruction} reveals that the 43 generators over the $4{,}080$ cyclic projective orbits in 16 variables span an explicit binary linear code that is strictly $16$-divisible ($4$-even). This code exhibits a large cyclic automorphism group and strict symplectic radical symmetries, offering a concrete blueprint for constructing new symmetric quantum CSS codes and quantum low-density parity-check (qLDPC) codes with natural transversal $T$-gates.

\subsection{Pseudorandom Sequences and Spread-Spectrum Communications}
Binary sequences with ideal periodic and aperiodic autocorrelation properties are critical in radar signal design, direct-sequence spread-spectrum communications (CDMA), and synchronization preambles in 5G/6G standards. Boolean functions invariant under cyclic group actions correspond directly to cyclically permuted sequences whose out-of-phase correlation values are governed by the Walsh-Hadamard transform.
Our Fourier-Walsh restriction principle provides an algebraic mechanism for decomposing character sums of cyclic sequences into fixed-point sub-sequences and uncoupled symplectic radical spaces. This enables new lower bounds on the cross-correlation of decimation sequences and offers an analytic filter for constructing sequence families with optimal peak-to-average power ratios (PAPR).

\subsection{Discrete Harmonic Analysis and Additive Combinatorics}
From the perspective of harmonic analysis on finite abelian groups, Lemma~\ref{lem:quadratic_reduction} constitutes a discrete restriction theorem for quadratic exponential sums: whenever a quadratic form is invariant under a group of odd-order automorphisms, its character sum factorizes cleanly across the Reynolds projection. In additive combinatorics, this principle provides a powerful method for computing the Gowers uniformity norms ($U^2$ and $U^3$) of sets with cyclic and affine symmetries, establishing that quadratic pseudorandomness on the global space $\F^n$ is entirely controlled by its restriction to the fixed-point subspace.

%======================================================================
\section{Conclusion and Open Problems}
\label{sec:conclusion}
%======================================================================

In this paper, we resolved the cubic St\u{a}nic\u{a}--Maitra conjecture for all even dimensions $n \not\equiv 0 \pmod{32}$. This result substantially expands beyond the previously known $n = 2p$ ($p$ prime) family of Cusick and Sanger~\cite{cusick2017}, providing uniform coverage for composite odd factors and power-of-two multiples $n \in \{2m, 4m, 8m, 16m\}$.

Two natural open problems remain:
\begin{enumerate}
    \item \textbf{The Structural Boundary at $t = 32$ and Breakdown of Direct Divisibility:} While the odd-order reduction framework (Lemma~\ref{lem:quadratic_reduction} and Theorem~\ref{thm:bent_restriction}) holds universally for all even dimensions $n$, the spectral obstruction via direct weight divisibility cannot be extended to $t = 32$. Indeed, direct divisibility already fails at $t = 32$: there exist explicit counterexamples of functions in the complete-orbit space $\mathcal{F}_{32}$ whose weights violate the 2-adic valuation bound $v_2(\wt(g)) \ge 16$ required to obstruct bentness, and the associated code is not 32-divisible. Thus, $t = 16$ represents the exact structural boundary of the divisible codes technique. Resolving the cubic conjecture for $n \equiv 0 \pmod{32}$ is not merely a pending computational verification, but an open theoretical problem requiring new algebraic mechanisms beyond direct linear code divisibility.
    \item \textbf{Degrees $d \ge 4$:} Extending the reduction framework to cubic or higher-degree derivatives requires non-quadratic exponential sum factorizations.
\end{enumerate}

%======================================================================
\begin{thebibliography}{99}
%======================================================================

\bibitem{calderon2024}
H.~Calder\'on-G\'omez, J.~Medina, and J.~Molina-Salazar,
``Short $k$-rotation symmetric Boolean functions,''
\emph{Discrete Applied Mathematics}, vol.~343, pp.~49--64, 2024.

\bibitem{cusick2017}
T.~W.~Cusick and E.~M.~Sanger,
``Rotation symmetric bent Boolean functions for $n = 2p$,''
arXiv preprint \href{https://arxiv.org/abs/1708.09313}{arXiv:1708.09313}, 2017.

\bibitem{meier1989}
W.~Meier and O.~Staffelbach,
``Nonlinearity criteria for cryptographic functions,''
in \emph{Proc. EUROCRYPT '89}, LNCS, vol.~434, pp.~549--562, Springer, 1989.

\bibitem{meng2010}
Q.~Meng, L.~S.~Chen, and F.-W.~Fu,
``On homogeneous rotation symmetric bent functions,''
\emph{Discrete Applied Mathematics}, vol.~158, no.~10, pp.~1111--1117, 2010.

\bibitem{pieprzyk1999}
J.~Pieprzyk and C.~X.~Qu,
``Fast hashing and rotation-symmetric functions,''
in \emph{Proc. Eurocrypt 1999}, LNCS, vol.~1592, pp.~95--105, Springer, 1999.

\bibitem{rothaus1976}
O.~S.~Rothaus,
``On `bent' functions,''
\emph{Journal of Combinatorial Theory, Series A}, vol.~20, no.~3, pp.~300--305, 1976.

\bibitem{stanica2008}
P.~St\u{a}nic\u{a} and S.~Maitra,
``Rotation symmetric Boolean functions---count and cryptographic properties,''
\emph{Discrete Applied Mathematics}, vol.~156, no.~10, pp.~1567--1580, 2008.

\bibitem{sun2026}
L.~Sun, Z.~Shi, J.~Liu, and F.-W.~Fu,
``On the conjecture about the nonexistence of homogeneous rotation symmetric bent functions,''
\emph{Designs, Codes and Cryptography}, vol.~94, no.~4, pp.~1023--1045, 2026.

\bibitem{ward1981}
H.~N.~Ward,
``Divisible codes,''
\emph{Archiv der Mathematik}, vol.~36, no.~1, pp.~485--494, 1981.

\bibitem{ward1990}
H.~N.~Ward,
``Weight polarization and divisibility,''
\emph{Discrete Mathematics}, vol.~83, no.~2--3, pp.~315--326, 1990.

\bibitem{ward2001}
H.~N.~Ward,
``Divisible codes---a survey,''
\emph{Serdica Mathematical Journal}, vol.~27, no.~4, pp.~263--278, 2001.

\bibitem{zhang2013}
G.~Zhang and X.~Gao,
``On the conjecture about the nonexistence of rotation symmetric bent functions,''
arXiv preprint \href{https://arxiv.org/abs/1303.2282}{arXiv:1303.2282}, 2013.

\end{thebibliography}

\end{document}
